\documentclass{article}
\usepackage[spanish]{babel}
\usepackage{graphicx,hyperref,listings,amsmath}
\title{Un art\'iculo de prueba}
\author{Ana P\'erez}
\date{2026}
\begin{document}
\maketitle
\begin{abstract}
Resumen breve con $x^2$.
\end{abstract}
\section{Introducci\'on}
Texto con \textbf{negrita}, \textit{it\'alica}, \emph{\'enfasis}, \texttt{c\'odigo} y
\underline{subrayado}. Marcador: fixture-tex-paper.
Ver \cite{knuth84,lamport94} y la ecuaci\'on~\eqref{eq:pit}.\footnote{Una nota al pie.}
\subsection{Listas}
\begin{itemize}
  \item uno
  \item dos
  \begin{enumerate}
    \item anidado
  \end{enumerate}
\end{itemize}
\begin{description}
  \item[T\'ermino] Su definici\'on.
\end{description}
\subsection{Matem\'aticas}
En l\'inea $E=mc^2$ y \(a+b\) y $\alpha_{i}~\text{texto}$.
\[ \int_0^1 x\,dx \]
\begin{equation}\label{eq:pit} a^2+b^2=c^2 \end{equation}
\begin{align} x &= 1 \\ y &= 2 \end{align}
\subsection{C\'odigo y citas}
\begin{lstlisting}[language=Python]
print("hola $no math$ \section{no}")
\end{lstlisting}
El comando \verb|\foo| en l\'inea.
\begin{quote}Una cita.\end{quote}
\href{https://example.org}{un enlace} y \url{https://example.org/b} -- rango --- raya ``comillas''.
\subsection{Tablas y figuras}
\begin{table}\centering
\begin{tabular}{|l|c|}\hline
A & B \\ \hline
1 & $x$ \\ \hline
\end{tabular}
\caption{Una tabla}
\end{table}
\begin{figure}\centering
\includegraphics[width=0.5\textwidth]{img/plot}
\caption{Un gr\'afico}
\end{figure}
\newpage
% un comentario que desaparece
\end{document}
